\chapter{Shortest Paths}\label{ch:shortest}

Finding the cheapest way from one place to another is the most common graph question in contests, and it appears in many
disguises: the fewest moves in a puzzle, the cheapest schedule, the latest start time consistent with a set of deadlines.
There is no single best algorithm. The right one depends on the edge weights (all equal, zero or one, non-negative,
arbitrary), on whether the graph has cycles, and on whether distances are needed from one source or between all pairs. This
chapter develops each algorithm from the single operation they share, relaxation, proves why it works, and shows the
modelling tricks that make it apply: states that remember extra information, reversed graphs, negated weights and systems
of inequalities.

\begin{chapterintro}{\silver\ \gold\ \plat}{Chapter~\ref{ch:graphs} (BFS and DFS) and Section~\ref{sec:heaps} (heaps).}
  \item choose the right shortest-path algorithm from the weights and the question;
  \item run Dijkstra's algorithm, prove it, and extend it to states, counts and path information;
  \item handle 0--1 weights with a deque, and negative weights with Bellman--Ford;
  \item find and print negative cycles, and detect unbounded longest paths;
  \item solve systems of difference constraints as shortest-path problems;
  \item compute all-pairs distances with Floyd--Warshall, and shortest and longest paths in DAGs.
\end{chapterintro}

% ======================================================================================================================
\section[Choosing the algorithm]{Choosing the algorithm \hfill \silver}\label{sec:sp-choose}

Figure~\ref{fig:spflow} is the decision procedure, and Table~\ref{tab:spalgos} the reference. Two questions settle almost
every case: what the weights look like, and whether the graph is a DAG.

\begin{figure}[htbp]
\centering
\begin{tikzpicture}[bus/.style={line width=0.8pt, draw=fgink!70}]
  \node[fq=3.6cm] (w) at (5.6, 0) {What are the edge weights?};
  \node[fend=1.9cm] (bfs) at (0, -2.0) {BFS};
  \node[fend=1.9cm] (zo) at (2.5, -2.0) {0--1 BFS};
  \node[fq=3.0cm] (nn) at (5.9, -2.0) {One source or all pairs?};
  \node[fq=3.0cm] (neg) at (11.55, -2.0) {Is the graph a DAG?};
  \node[fend=2.2cm] (dij) at (4.6, -4.3) {Dijkstra};
  \node[fend=2.2cm] (fw) at (7.2, -4.3) {Floyd--Warshall ($n \le 500$)};
  \node[fend=2.2cm] (dag) at (10.2, -4.3) {DP in topological order};
  \node[fend=2.2cm] (bf) at (12.9, -4.3) {Bellman--Ford};
  \draw[bus] (w.south) -- (5.6, -0.95) (0, -0.95) -- (11.55, -0.95);
  \draw[farr] (0, -0.95) -- node[fyes, pos=0.4] {all equal} (bfs.north);
  \draw[farr] (2.5, -0.95) -- node[fyes, pos=0.4] {0 or 1} (zo.north);
  \draw[farr] (5.9, -0.95) -- node[fyes, pos=0.4] {$\ge 0$} (nn.north);
  \draw[farr] (11.55, -0.95) -- node[fno, pos=0.4] {some $< 0$} (neg.north);
  \draw[bus] (nn.south) -- (5.9, -3.2) (4.6, -3.2) -- (7.2, -3.2);
  \draw[farr] (4.6, -3.2) -- node[fyes, pos=0.35] {one} (dij.north);
  \draw[farr] (7.2, -3.2) -- node[fyes, pos=0.35] {all pairs} (fw.north);
  \draw[bus] (neg.south) -- (11.55, -3.2) (10.2, -3.2) -- (12.9, -3.2);
  \draw[farr] (10.2, -3.2) -- node[fyes, pos=0.35] {yes} (dag.north);
  \draw[farr] (12.9, -3.2) -- node[fno, pos=0.35] {no} (bf.north);
\end{tikzpicture}
\caption{Which shortest-path algorithm? A DAG allows any weights (and even longest paths) with a simple DP; negative weights
in a graph with cycles need Bellman--Ford, or Floyd--Warshall for all pairs.}
\label{fig:spflow}
\end{figure}

\begin{table}[htbp]
\centering
\small
\begin{tabular}{@{}P{3.9cm} P{3.3cm} l P{3.8cm}@{}}
  \toprule
  \thd{Weights and question} & \thd{Algorithm} & \thd{Time} & \thd{Remarks} \\
  \midrule
  all equal & BFS & $\Oh{n + m}$ & Chapter~\ref{ch:graphs} \\
  0 or 1 & 0--1 BFS & $\Oh{n + m}$ & deque instead of queue \\
  non-negative, one source & Dijkstra & $\Oh{(n + m)\log n}$ & the workhorse \\
  any, one source & Bellman--Ford & $\Oh{nm}$ & detects negative cycles \\
  any, all pairs & Floyd--Warshall & $\Oh{n^3}$ & three lines; $n \le 500$ \\
  any, DAG & DP in topological order & $\Oh{n + m}$ & longest paths too \\
  non-negative, all pairs, sparse & $n \times$ Dijkstra & $\Oh{n (n + m)\log n}$ & Johnson if some are negative \\
  \bottomrule
\end{tabular}
\caption{Shortest-path algorithms at a glance.}
\label{tab:spalgos}
\end{table}

% ======================================================================================================================
\section[Relaxation]{Relaxation and shortest-path trees \hfill \silver}\label{sec:relax}

Every algorithm in this chapter maintains tentative distances $d(v)$, upper bounds on the true distances $\delta(v)$, with
$d(s) = 0$ and $d(v) = \infty$ elsewhere at the start. The only way they ever change is by \defn{relaxation}: for an edge
$u \to v$ of weight $w$,
\[
  \text{if } d(u) + w < d(v) \text{ then } d(v) \gets d(u) + w, \ \ \mathit{par}(v) \gets u .
\]
Relaxation keeps every $d(v)$ an upper bound, because $d(u) + w$ is the length of a real walk to $v$. At the end, when no
edge can be relaxed, every edge satisfies the \emph{triangle inequality} $d(v) \le d(u) + w$, and the distances are exact:
along any shortest path $s = v_0, v_1, \ldots, v_k$, induction gives $d(v_i) \le \delta(v_i)$. The algorithms differ only in
the \emph{order} in which they relax edges, and that order is what they need to be both correct and fast.

The parent pointers form a \defn{shortest-path tree}: following $\mathit{par}$ from any vertex back to $s$ spells out a
shortest path, in reverse. So storing one parent per vertex is all it takes to print paths, and it works because shortest
paths have \emph{optimal substructure}: every prefix of a shortest path is itself a shortest path.

% ======================================================================================================================
\section[Dijkstra's algorithm]{Dijkstra's algorithm \hfill \gold}\label{sec:dijkstra}

\defn{Dijkstra's algorithm} handles non-negative weights. It repeatedly takes the vertex with the smallest tentative distance
that is not yet \emph{final}, declares it final, and relaxes its outgoing edges (Figure~\ref{fig:dijkstra}). A priority queue
of (distance, vertex) pairs finds that vertex; instead of decreasing keys, the implementation pushes a new pair whenever a
distance improves and skips popped pairs whose distance is out of date (\emph{lazy deletion}).

\begin{figure}[htbp]
\centering
\newcommand{\djpanel}[2]{%
\begin{tikzpicture}[x=0.86cm, y=0.86cm]
  \coordinate (S) at (0, 1); \coordinate (A) at (1.3, 2); \coordinate (B) at (1.3, 0);
  \coordinate (C) at (2.8, 2); \coordinate (D) at (2.8, 0); \coordinate (T) at (4.1, 1);
  \foreach \ea/\eb in {S/A, S/B, A/B, A/C, B/C, B/D, C/T, D/T} \draw[ed] (\ea) -- (\eb);
  \foreach \ea/\eb in {#2} \draw[bed] (\ea) -- (\eb);
  \foreach \ea/\eb/\w in {S/A/2, S/B/7, A/B/3, A/C/4, B/C/1, B/D/4, C/T/5, D/T/1}
    \node[wl, font=\sffamily\tiny] at ($(\ea)!0.5!(\eb)$) {\w};
  \foreach \vn/\sy/\dv in {#1} \node[vtxs, \sy] at (\vn) {\dv};
  \foreach \vn/\pos in {S/left, A/above, B/below, C/above, D/below, T/right}
    \node[font=\sffamily\tiny\bfseries, text=fggray, \pos=5pt] at (\vn) {\vn};
\end{tikzpicture}}
\begin{filmstrip}[3]
  \begin{stepbox}{1. start: $d(S) = 0$}
    \djpanel{S/cool/0, A/dim/$\infty$, B/dim/$\infty$, C/dim/$\infty$, D/dim/$\infty$, T/dim/$\infty$}{S/S}
  \end{stepbox}
  \begin{stepbox}{2. finalise S (0), relax}
    \djpanel{S/hot/0, A/cool/2, B/cool/7, C/dim/$\infty$, D/dim/$\infty$, T/dim/$\infty$}{S/A, S/B}
  \end{stepbox}
  \begin{stepbox}{3. finalise A (2): B improves to 5}
    \djpanel{S/good/0, A/hot/2, B/cool/5, C/cool/6, D/dim/$\infty$, T/dim/$\infty$}{S/A, A/B, A/C}
  \end{stepbox}
  \begin{stepbox}{4. finalise B (5)}
    \djpanel{S/good/0, A/good/2, B/hot/5, C/cool/6, D/cool/9, T/dim/$\infty$}{S/A, A/B, A/C, B/D}
  \end{stepbox}
  \begin{stepbox}{5. finalise C (6): T gets 11}
    \djpanel{S/good/0, A/good/2, B/good/5, C/hot/6, D/cool/9, T/cool/11}{S/A, A/B, A/C, B/D, C/T}
  \end{stepbox}
  \begin{stepbox}{6. finalise D (9): T improves to 10}
    \djpanel{S/good/0, A/good/2, B/good/5, C/good/6, D/hot/9, T/cool/10}{S/A, A/B, A/C, B/D, D/T}
  \end{stepbox}
\end{filmstrip}
\caption{Dijkstra's algorithm from S. Green vertices are final, red is the vertex being finalised, blue ones are in the
queue with a tentative distance, and thick blue edges are the current parent pointers. The last step finalises T with
distance 10, and the parent pointers form the shortest-path tree.}
\label{fig:dijkstra}
\end{figure}

\begin{pseudo}{Dijkstra}
\begin{algorithmic}[1]
\State $d(s) \gets 0$, $d(v) \gets \infty$ for $v \ne s$; push $(0, s)$
\While{the priority queue is not empty}
  \State pop the pair $(x, u)$ with the smallest $x$
  \If{$x \ne d(u)$} \textbf{continue} \Comment{an outdated pair}
  \EndIf
  \ForAll{edges $u \to v$ of weight $w$}
    \If{$d(u) + w < d(v)$} \State $d(v) \gets d(u) + w$; push $(d(v), v)$ \EndIf
  \EndFor
\EndWhile
\end{algorithmic}
\end{pseudo}

\emph{Why it is correct.} Claim: when $u$ is popped with $x = d(u)$, $d(u) = \delta(u)$. Suppose not, and let $u$ be the first
vertex finalised with $d(u) > \delta(u)$. Take a shortest path from $s$ to $u$, and let $y$ be its first vertex that is not
final, with predecessor $x'$ on the path (which is final, hence correct). When $x'$ was finalised, the edge $x' \to y$ was
relaxed, so $d(y) = \delta(y)$. Since weights are non-negative and $y$ lies on the path to $u$, $\delta(y) \le \delta(u) <
d(u)$, so $d(y) < d(u)$, and the queue would have popped $y$ before $u$: a contradiction. This is exactly where
non-negativity is used, and why a single negative edge can break the algorithm.

\emph{Complexity.} Every edge pushes at most one pair, so the queue sees at most $m + 1$ pushes and pops: $\Oh{(n + m)\log m}
= \Oh{(n + m) \log n}$. On dense graphs ($m \approx n^2$) the $\Oh{n^2}$ version without a heap, which scans all vertices for
the minimum, is faster.

\begin{code}
vector<ll> dist(n, LLONG_MAX);
minpq<pll> pq;  // (distance, vertex)
dist[src] = 0;
pq.push({0, src});
while (!pq.empty()) {
    auto [d, u] = pq.top();
    pq.pop();
    if (d != dist[u]) continue;  // an outdated entry ("lazy deletion")
    for (auto [v, w] : adj[u])
        if (d + w < dist[v]) dist[v] = d + w, pq.push({dist[v], v});
}
\end{code}

\usetemplate{Dijkstra}{Ctrl+Alt+T, J}{cpdijkstra}

The complete program for the basic problem is the worked problem Shortest Routes I (Chapter~\ref{ch:worked}). The two problems
below show the two most important extensions: richer states, and information gathered along the shortest paths.

\subsection{Dijkstra on states}

\begin{problem}{Flight Discount (CSES)}
Find the cheapest route from city 1 to city $n$ ($n \le 10^5$, $m \le 2\cdot10^5$ one-way flights) when exactly one flight may
be taken at half price (rounded down).
\solution The cost of the rest of the trip depends on the city \emph{and} on whether the coupon is still available. So the
search runs on the graph of states $(v, \mathit{used})$, which has $2n$ vertices: a flight $u \to v$ of price $w$ becomes the
edges $(u, 0) \to (v, 0)$ and $(u, 1) \to (v, 1)$ of weight $w$, plus the coupon edge $(u, 0) \to (v, 1)$ of weight
$\lfloor w/2 \rfloor$. One Dijkstra on this layered graph answers the question.
\end{problem}

\begin{code}
void solve() {
    int n, m;
    cin >> n >> m;
    vector<vector<pll>> adj(n + 1);
    rep(i, m) {
        ll a, b, c;
        cin >> a >> b >> c;
        adj[a].push_back({b, c});
    }
    // state (v, used): at city v, the coupon already used (1) or not (0)
    vector<array<ll, 2>> dist(n + 1, {INF, INF});
    minpq<tuple<ll, int, int>> pq;  // (distance, city, used)
    dist[1][0] = 0;
    pq.push({0, 1, 0});
    while (!pq.empty()) {
        auto [d, u, used] = pq.top();
        pq.pop();
        if (d != dist[u][used]) continue;  // an outdated entry
        for (auto [v, w] : adj[u]) {
            if (d + w < dist[v][used]) {  // pay the full price
                dist[v][used] = d + w;
                pq.push({dist[v][used], v, used});
            }
            if (!used && d + w / 2 < dist[v][1]) {  // use the coupon now
                dist[v][1] = d + w / 2;
                pq.push({dist[v][1], v, 1});
            }
        }
    }
    print(min(dist[n][0], dist[n][1]));
}
\end{code}

\begin{trick}[Layered graphs]
``At most $k$ free edges'', ``at most $k$ stops'', ``the parity of the number of steps'', ``which key you carry'': whenever the
future depends on a small amount of extra information, make it part of the vertex. The graph grows by that factor ($k + 1$
layers, 2 parities, $2^{\text{keys}}$ key sets), and the algorithm stays the same. The alternative for Flight Discount, which
avoids the layered graph, is also worth knowing: run Dijkstra from 1 and, on the reversed graph, from $n$; the answer is the
minimum over all flights $u \to v$ of $d_1(u) + \lfloor w/2\rfloor + d_n(v)$.
\end{trick}

\subsection{Counting and comparing shortest paths}

\begin{problem}{Investigation (CSES)}
For the cheapest routes from city 1 to city $n$, find their price, how many there are (modulo $10^9 + 7$), and the fewest and
the most flights among them. All prices are positive.
\solution Dijkstra finalises vertices in order of distance, and with positive weights every predecessor of $v$ on a cheapest
route is finalised before $v$. So when a vertex $u$ is finalised its counts are complete and can be pushed along its edges:
a strictly cheaper route to $v$ overwrites $v$'s information, an equally cheap one merges into it.
\end{problem}

\begin{code}
void solve() {
    int n, m;
    cin >> n >> m;
    vector<vector<pll>> adj(n + 1);
    rep(i, m) {
        ll a, b, c;
        cin >> a >> b >> c;
        adj[a].push_back({b, c});
    }
    vector<ll> dist(n + 1, LLONG_MAX), ways(n + 1, 0);  // ways mod MOD
    vi fewest(n + 1, 0), most(n + 1, 0);  // flights on the cheapest routes
    minpq<pll> pq;
    dist[1] = 0, ways[1] = 1;
    pq.push({0, 1});
    while (!pq.empty()) {
        auto [d, u] = pq.top();
        pq.pop();
        if (d != dist[u]) continue;
        // u is final: its counts are complete, so push them forward
        for (auto [v, w] : adj[u]) {
            if (d + w < dist[v]) {  // a strictly cheaper route: start over
                dist[v] = d + w;
                ways[v] = ways[u];
                fewest[v] = fewest[u] + 1, most[v] = most[u] + 1;
                pq.push({dist[v], v});
            } else if (d + w == dist[v]) {  // another cheapest route
                ways[v] = (ways[v] + ways[u]) % MOD;
                fewest[v] = min(fewest[v], fewest[u] + 1);
                most[v] = max(most[v], most[u] + 1);
            }
        }
    }
    print(dist[n], ways[n], fewest[n], most[n]);
}
\end{code}

\begin{trick}[More Dijkstra variations]
\begin{itemize}
  \item \textbf{Which edges lie on some shortest path?} Run Dijkstra from $s$ and, on the reversed graph, from $t$; the edge
        $u \to v$ of weight $w$ lies on a shortest path exactly when $d_s(u) + w + d_t(v) = d_s(t)$.
  \item \textbf{Minimax paths} (minimise the largest edge on the path): replace $d(u) + w$ by $\max(d(u), w)$; the same proof
        works. Kruskal gives the same answers for all pairs at once (Chapter~\ref{ch:dsu}).
  \item \textbf{The $k$ shortest walks}: let every vertex be popped up to $k$ times, keeping a counter per vertex; the $k$-th
        pop of $t$ is the $k$-th shortest walk (CSES Flight Routes).
  \item \textbf{Multiple sources}: push them all with distance 0, exactly as in multi-source BFS.
  \item \textbf{Distances to one target}: run Dijkstra from the target on the reversed graph.
\end{itemize}
\end{trick}

\begin{pitfall}[The classic Dijkstra bugs]
Forgetting the \cc{if (d != dist[u]) continue;} line keeps the answer correct but can make the algorithm quadratic.
Declaring \cc{priority_queue<pll>} without \cc{greater} (or \cc{minpq}) turns it into a max-heap, which silently explores in
the wrong order. Using \cc{int} distances overflows with $10^5$ edges of weight $10^9$. And with a negative edge, the algorithm
is simply wrong, even when there is no negative cycle.
\end{pitfall}

% ======================================================================================================================
\section[0--1 BFS]{0--1 BFS \hfill \gold}\label{sec:01bfs}

When every weight is 0 or 1, Dijkstra's priority queue can be replaced by a deque. Relaxing along a 0-edge produces a vertex
with the \emph{same} distance as the current one, so it goes to the front; a 1-edge produces the next distance, so it goes to the
back. The deque then always holds tentative distances $d$ at the front and $d + 1$ at the back (Figure~\ref{fig:01bfs}), which is
exactly the order Dijkstra would use, at $\Oh{1}$ per operation: $\Oh{n + m}$ in total.

\begin{figure}[htbp]
\centering
\begin{tikzpicture}
  \arrs{Q}{0}{0}{4/cool, 4/cool, 4/cool, 5/warn, 5/warn}
  \node[blbl, anchor=east] at (-0.3, 0) {deque};
  \node[lbl, anchor=south] at ($(Q-1.north)+(0,0.05)$) {front};
  \node[lbl, anchor=south] at ($(Q-5.north)+(0,0.05)$) {back};
  \draw[->, line width=1pt, draw=fgblue] ($(Q-1.west)+(-0.9, -0.75)$) node[below, font=\sffamily\scriptsize,
    text=fgblue] {0-edge: push front, distance $d$} to[out=0, in=-120] ($(Q-1.south west)+(0.15, 0)$);
  \draw[->, line width=1pt, draw=fgorange!85!black] ($(Q-5.east)+(0.9, -0.75)$) node[below, font=\sffamily\scriptsize,
    text=fgorange!80!black] {1-edge: push back, distance $d + 1$} to[out=180, in=-60] ($(Q-5.south east)+(-0.15, 0)$);
\end{tikzpicture}
\caption{0--1 BFS keeps the deque sorted by distance: vertices at distance $d$ (blue) in front, $d + 1$ (orange) behind.}
\label{fig:01bfs}
\end{figure}

\begin{problem}{Fewest walls to break}
In an $n \times m$ grid of free cells and walls, moving into a free cell costs nothing and moving into a wall cell (breaking it)
costs 1. Find the minimum cost of a path from the top-left to the bottom-right cell (counting the start cell if it is a wall).
\solution A 0--1 BFS over cells, where the weight of a move is 1 exactly when the target cell is a wall.
\end{problem}

\begin{code}
void solve() {
    int n, m;
    cin >> n >> m;
    vector<string> g(n);
    cin >> g;
    const int dr[4] = {-1, 1, 0, 0}, dc[4] = {0, 0, -1, 1};
    // entering a wall cell costs 1 (break it), a free cell costs 0
    vvi dist(n, vi(m, INF));
    deque<pii> dq;
    dist[0][0] = g[0][0] == '#';
    dq.push_back({0, 0});
    while (!dq.empty()) {
        auto [r, c] = dq.front();
        dq.pop_front();
        rep(d, 4) {
            int nr = r + dr[d], nc = c + dc[d];
            if (nr < 0 || nr >= n || nc < 0 || nc >= m) continue;
            int w = g[nr][nc] == '#';
            if (dist[r][c] + w >= dist[nr][nc]) continue;
            dist[nr][nc] = dist[r][c] + w;
            if (w == 0) dq.push_front({nr, nc});  // same distance: front
            else dq.push_back({nr, nc});  // one more: back
        }
    }
    print(dist[n - 1][m - 1]);
}
\end{code}

\usetemplate{0-1 BFS}{Ctrl+Alt+T, 0}{cp01bfs}

The same idea extends to small integer weights: with weights in $\{0, 1, \ldots, k\}$, keep $k + 1$ buckets (a circular array
of queues indexed by distance modulo $k + 1$), known as \emph{Dial's algorithm}, for $\Oh{n k + m}$.

% ======================================================================================================================
\section[Negative weights: Bellman--Ford]{Negative weights: Bellman--Ford \hfill \plat}\label{sec:bellman}

With negative edges, the order of relaxations can no longer be chosen greedily. \defn{Bellman--Ford} does not choose at all:
it relaxes \emph{every} edge, in any order, $n - 1$ times. After round $k$, every vertex whose shortest path uses at most $k$
edges has its correct distance (induction on $k$: the last edge of such a path is relaxed in round $k$, after its start became
correct). If there is no negative cycle reachable from the source, shortest paths are simple and have at most $n - 1$ edges,
so $n - 1$ rounds suffice. If an edge can still be relaxed in round $n$, there is a negative cycle: distances would keep falling
forever.

\begin{problem}{Cycle Finding (CSES)}
Determine whether a directed graph ($n \le 2500$, $m \le 5000$) contains a negative cycle, and if so, print one.
\solution Start every distance at 0, as if a virtual source had 0-edges to every vertex; then every negative cycle is
reachable. Run $n$ rounds and remember a vertex relaxed in the last one. That vertex may hang off the cycle rather than lie on it
(Figure~\ref{fig:negcycle}), but its parent pointers lead back into the cycle: after $n$ steps back along them we are certainly on
it, and one more walk around recovers the cycle.
\end{problem}

\begin{figure}[htbp]
\centering
\begin{tikzpicture}
  \node[vtx] (v1) at (0, 0) {1};
  \node[vtx] (v2) at (2, 0) {2};
  \node[vtx] (v3) at (3.4, 1.1) {3};
  \node[vtx] (v4) at (3.4, -1.1) {4};
  \node[vtx, hot] (v5) at (5.6, -1.1) {5};
  \draw[ded] (v1) -- node[wl] {2} (v2);
  \draw[hded] (v2) -- node[wl] {$-3$} (v3);
  \draw[hded] (v3) -- node[wl] {1} (v4);
  \draw[hded] (v4) -- node[wl] {$-1$} (v2);
  \draw[ded] (v4) -- node[wl] {2} (v5);
  \node[notel, anchor=west, text width=5.4cm] at (6.3, 0) {The cycle $2 \to 3 \to 4 \to 2$ has weight $-3$. A vertex relaxed
    in round $n$ can be off the cycle, like 5; walking back along parent pointers ($5 \to 4 \to 3 \to 2 \to 4 \to \cdots$)
    enters the cycle within $n$ steps.};
\end{tikzpicture}
\caption{A negative cycle (red) and a vertex that depends on it. Everything reachable from a negative cycle has distance
$-\infty$.}
\label{fig:negcycle}
\end{figure}

\begin{code}
void solve() {
    int n, m;
    cin >> n >> m;
    vector<array<ll, 3>> e(m);
    for (auto& [a, b, c] : e) cin >> a >> b >> c;
    // all distances start at 0: a virtual source reaches every vertex
    vector<ll> dist(n + 1, 0);
    vi par(n + 1, 0);
    int last = 0;  // a vertex relaxed in round n, if any
    rep(round, n) {
        last = 0;
        for (auto [a, b, c] : e)
            if (dist[a] + c < dist[b])
                dist[b] = dist[a] + c, par[b] = a, last = b;
    }
    if (!last) {
        print("NO");
        return;
    }
    // last may hang off the cycle: n steps back along par land on it
    int x = last;
    rep(i, n) x = par[x];
    vi cycle{x};
    for (int y = par[x]; y != x; y = par[y]) cycle.push_back(y);
    cycle.push_back(x);
    reverse(all(cycle));
    print("YES");
    print(cycle);
}
\end{code}

\usetemplate{Bellman-Ford (negative edges)}{Ctrl+Alt+T, Shift+W}{cpbellman}

\subsection{Longest paths and unbounded answers}

Negating all weights turns a longest-path problem into a shortest-path problem, and a positive cycle into a negative one. The
subtle part is deciding when the answer is unbounded: only if a cycle that can be improved forever is both reachable from the
start \emph{and} able to reach the target.

\begin{problem}{High Score (CSES)}
Each tunnel $a \to b$ changes your score by $x$ (possibly negative). Find the maximum score of a route from room 1 to room $n$,
or $-1$ if it can be arbitrarily large.
\solution Run Bellman--Ford for $n - 1$ rounds on the negated scores. Any vertex that can still be improved afterwards is on or
after a positive-score cycle reachable from 1. Then run $n$ more rounds that mark every such vertex and spread the mark along
edges; if room $n$ gets marked, the answer is unbounded.
\end{problem}

\begin{code}
void solve() {
    int n, m;
    cin >> n >> m;
    vector<array<ll, 3>> e(m);
    for (auto& [a, b, c] : e) cin >> a >> b >> c;
    // longest path = shortest path with negated scores
    const ll UNREACHED = LLONG_MAX;
    vector<ll> dist(n + 1, UNREACHED);
    dist[1] = 0;
    rep(round, n - 1)
        for (auto [a, b, c] : e)
            if (dist[a] != UNREACHED && dist[a] - c < dist[b])
                dist[b] = dist[a] - c;
    // anything still improvable lies on or after a positive-score cycle;
    // n more rounds spread that mark to everything reachable from it
    vector<bool> inf(n + 1, false);
    rep(round, n)
        for (auto [a, b, c] : e) {
            if (dist[a] == UNREACHED) continue;
            if (inf[a] || dist[a] - c < dist[b])
                inf[b] = true, dist[b] = dist[a] - c;
        }
    print(inf[n] ? -1 : -dist[n]);
}
\end{code}

\begin{fact}[SPFA]
The ``Shortest Path Faster Algorithm'' is Bellman--Ford with a queue of vertices whose distance changed. It is often much
faster in practice, but its worst case is still $\Oh{nm}$, and problem setters build that case on purpose. Use it where
Bellman--Ford is intended anyway, such as inside min-cost flow with negative costs, never as a replacement for Dijkstra.
\end{fact}

\begin{trick}[Johnson's reweighting]
With negative edges but no negative cycle, all-pairs distances in a sparse graph are best computed with \emph{potentials}: run
Bellman--Ford once from a virtual source to get $h(v)$, then run Dijkstra from every vertex with the weights
$w'(u, v) = w(u, v) + h(u) - h(v)$, which are non-negative by the triangle inequality. Every path from $s$ to $t$ changes by the
same amount $h(s) - h(t)$, so shortest paths are preserved and $\delta(s, t) = \delta'(s, t) - h(s) + h(t)$. Min-cost flow uses
exactly these potentials to keep running Dijkstra after negative residual edges appear (Chapter~\ref{ch:flows}).
\end{trick}

% ======================================================================================================================
\section[Difference constraints]{Systems of difference constraints \hfill \plat}\label{sec:diffcons}

A surprising number of scheduling problems reduce to a system of inequalities of the form $x_b - x_a \le c$: ``task $b$ starts
at most $c$ after task $a$'', ``the prefix sum up to $j$ exceeds the one up to $i$ by at least 3''. Such a \defn{system of
difference constraints} is exactly the triangle inequality of shortest paths. Make a vertex per variable and an edge
$a \to b$ of weight $c$ per constraint, and add a virtual source with 0-edges to all vertices. Then:
\begin{itemize}
  \item if the graph has a negative cycle, the system has no solution: adding up the constraints around the cycle gives
        $0 \le (\text{negative number})$;
  \item otherwise the shortest distances $x_v = d(v)$ satisfy every constraint, because $d(b) \le d(a) + c$ for every edge.
\end{itemize}
So Bellman--Ford decides feasibility and constructs a solution (Figure~\ref{fig:diffcons}). Among all solutions with every
$x_v \le 0$, it even finds the componentwise largest one. Equalities are two inequalities, and $x_b - x_a \ge c$ is
$x_a - x_b \le -c$.

\begin{figure}[htbp]
\centering
\begin{tikzpicture}
  \node[vtx, dim] (s) at (-2.4, 0) {s};
  \node[vtx] (x1) at (0, 1.3) {$x_1$};
  \node[vtx] (x2) at (2.4, 0) {$x_2$};
  \node[vtx] (x3) at (0, -1.3) {$x_3$};
  \foreach \tg in {x1, x2, x3} \draw[->, line width=0.6pt, draw=fgline, dashed] (s) to[bend left=8] node[font=\tiny,
    text=fggray, fill=white, inner sep=0.5pt] {0} (\tg);
  \draw[ded] (x1) -- node[wl] {3} (x2);
  \draw[ded] (x2) -- node[wl] {$-2$} (x3);
  \draw[ded] (x3) -- node[wl] {1} (x1);
  \node[font=\sffamily\scriptsize\bfseries, text=fgblue] at (0.55, 1.85) {$-1$};
  \node[font=\sffamily\scriptsize\bfseries, text=fgblue] at (3.0, 0.4) {0};
  \node[font=\sffamily\scriptsize\bfseries, text=fgblue] at (0.55, -1.85) {$-2$};
  \node[notel, anchor=west, text width=5.8cm] at (3.7, 0) {Constraints $x_2 - x_1 \le 3$, $x_3 - x_2 \le -2$, $x_1 - x_3 \le
    1$. The cycle $x_1 \to x_2 \to x_3 \to x_1$ has weight $2 \ge 0$, so the system is feasible, and the distances from the
    virtual source $s$ (blue) are a solution: $x = (-1, 0, -2)$.};
\end{tikzpicture}
\caption{A system of difference constraints as a graph. Each constraint $x_b - x_a \le c$ is an edge $a \to b$ of weight $c$;
shortest distances from a virtual source satisfy all constraints at once.}
\label{fig:diffcons}
\end{figure}

\begin{problem}{Difference constraints}
Given $n$ variables and $m$ constraints $x_b - x_a \le c$, print integer values satisfying all of them, or \texttt{NO}.
\end{problem}

\begin{code}
void solve() {
    int n, m;
    cin >> n >> m;
    // constraint x_b - x_a <= c becomes the edge a -> b of weight c
    vector<array<ll, 3>> e(m);
    for (auto& [a, b, c] : e) cin >> a >> b >> c;
    vector<ll> x(n + 1, 0);  // a virtual source with 0-edges to everyone
    rep(round, n + 1) {
        bool changed = false;
        for (auto [a, b, c] : e)
            if (x[a] + c < x[b]) x[b] = x[a] + c, changed = true;
        if (!changed) {  // shortest distances satisfy every constraint
            print(vi(x.begin() + 1, x.end()));
            return;
        }
    }
    print("NO");  // still changing after n + 1 rounds: a negative cycle
}
\end{code}

\usetemplate{System of difference constraints}{Ctrl+Alt+U, Shift+I}{cpdiffcons}

\begin{recognize}
Variables that are only compared by differences: start times with precedence and deadline constraints, positions on a line
with minimum and maximum gaps, prefix sums with bounds on segment sums (``every window of length $k$ contains at least $a$
chosen elements''). When the constraints involve a parameter, binary search the parameter and test feasibility with
Bellman--Ford.
\end{recognize}

% ======================================================================================================================
\section[All pairs: Floyd--Warshall]{All pairs: Floyd--Warshall \hfill \gold}\label{sec:floyd}

\defn{Floyd--Warshall} computes all $n^2$ distances in $\Oh{n^3}$ with three nested loops. It is a dynamic programme over the set
of vertices allowed as \emph{intermediate} stops: after the iteration for $k$, \cc{d[i][j]} is the length of the shortest path
from $i$ to $j$ whose intermediate vertices all lie in $\{1, \ldots, k\}$. Allowing $k$ as a new stop either does not help, or
the best path passes through $k$ exactly once, splitting into a path $i \leadsto k$ and a path $k \leadsto j$ that use only
earlier stops (Figure~\ref{fig:floyd}):
\[
  d_k[i][j] = \min\bigl(d_{k-1}[i][j],\ d_{k-1}[i][k] + d_{k-1}[k][j]\bigr) .
\]
The update can be done in place, because row $k$ and column $k$ do not change during iteration $k$. It handles negative edges,
and afterwards $d[i][i] < 0$ exactly when $i$ lies on a negative cycle.

\begin{figure}[htbp]
\centering
\begin{tikzpicture}
  \node[vtx] (i) at (0, 0) {$i$};
  \node[vtx, hot] (k) at (3, 1.4) {$k$};
  \node[vtx] (j) at (6, 0) {$j$};
  \draw[line width=1pt, draw=fgblue, decorate, decoration={snake, amplitude=1.2pt, segment length=7pt, post length=2mm}, ->] (i) --
    node[above left, font=\sffamily\scriptsize, text=fgblue] {stops $< k$} (k);
  \draw[line width=1pt, draw=fgblue, decorate, decoration={snake, amplitude=1.2pt, segment length=7pt, post length=2mm}, ->] (k) --
    node[above right, font=\sffamily\scriptsize, text=fgblue] {stops $< k$} (j);
  \draw[line width=1pt, draw=fggray, decorate, decoration={snake, amplitude=1.2pt, segment length=7pt, post length=2mm}, ->] (i) --
    node[below, font=\sffamily\scriptsize, text=fggray] {the best path with stops $< k$ only} (j);
\end{tikzpicture}
\caption{The Floyd--Warshall recurrence. With stops up to $k$ allowed, the best path from $i$ to $j$ either avoids $k$ (grey)
or visits it once, and then both halves use only stops below $k$ (blue).}
\label{fig:floyd}
\end{figure}

\begin{problem}{Shortest Routes II (CSES)}
Answer $10^5$ queries for the shortest distance between two of $n \le 500$ cities connected by two-way roads.
\end{problem}

\begin{code}
void solve() {
    int n, m, q;
    cin >> n >> m >> q;
    const ll NONE = LLONG_MAX / 4;  // large, but NONE + NONE does not overflow
    vector<vector<ll>> d(n + 1, vector<ll>(n + 1, NONE));
    rep(v, 1, n + 1) d[v][v] = 0;
    rep(i, m) {
        ll a, b, c;
        cin >> a >> b >> c;
        chmin(d[a][b], c);  // keep the cheapest of parallel roads
        chmin(d[b][a], c);
    }
    rep(k, 1, n + 1) rep(i, 1, n + 1) rep(j, 1, n + 1)  // k outermost!
        chmin(d[i][j], d[i][k] + d[k][j]);
    while (q--) {
        int a, b;
        cin >> a >> b;
        print(d[a][b] == NONE ? -1 : d[a][b]);
    }
}
\end{code}

\usetemplate{Floyd-Warshall (all pairs)}{Ctrl+Alt+T, W}{cpfloyd}

\begin{trick}[Floyd--Warshall is a template for more]
Replacing $(\min, +)$ by (\emph{or}, \emph{and}) computes the transitive closure, which with bitsets runs in $\Oh{n^3/64}$
(Section~\ref{sec:bitset}); $(\max, \min)$ computes widest paths. Because the outer loop adds vertices one at a time in any
order you like, processing them sorted by some value answers ``the shortest path using only vertices of value at most $x$''
for every $x$ in one run. And adding a single new edge $(a, b)$ to a finished distance matrix costs only $\Oh{n^2}$:
$d[i][j] \gets \min(d[i][j], d[i][a] + w + d[b][j])$.
\end{trick}

\begin{pitfall}[Loop order and infinity]
The loop over $k$ must be the outermost; with any other order the result is wrong on some inputs. And ``infinity'' must survive
additions: with \cc{LLONG_MAX}, \cc{d[i][k] + d[k][j]} overflows, so use a quarter of it (or skip pairs where either term is
infinite, as the template does).
\end{pitfall}

% ======================================================================================================================
\section[Shortest and longest paths in DAGs]{Shortest and longest paths in DAGs \hfill \gold}\label{sec:dagpaths}

In a DAG, process the vertices in topological order (Chapter~\ref{ch:dsu}): when a vertex is reached, every edge into it has
already been relaxed, so its distance is final, and relaxing its outgoing edges once is enough. This takes $\Oh{n + m}$ and
allows any weights. Most importantly, it solves \emph{longest} paths, which are hopeless in general graphs (they generalise
Hamiltonian paths) but easy in DAGs: just maximise instead of minimise (Figure~\ref{fig:dagpath}). Every DP is, at heart, a
shortest or longest path in the DAG of its states (Part~V).

\begin{figure}[htbp]
\centering
\begin{tikzpicture}[x=1.6cm]
  \foreach \vn/\xx/\yy/\val in {1/0/0/1, 2/1/0.8/2, 3/1/-0.8/2, 4/2/0.8/3, 5/2/-0.8/3, 6/3/0/4} {
    \node[vtx] (n\vn) at (\xx, \yy) {\vn};
    \node[font=\sffamily\scriptsize\bfseries, text=fgblue, above=1pt] at (n\vn.north) {\val};
  }
  \draw[ded] (n1) -- (n2); \draw[ded] (n1) -- (n3); \draw[ded] (n3) -- (n4); \draw[ded] (n3) -- (n5);
  \draw[ded] (n2) -- (n5); \draw[ded] (n5) -- (n6);
  \draw[hded] (n2) -- (n4); \draw[hded] (n4) -- (n6);
  \draw[hded] (n1) -- (n2);
  \node[notel, anchor=west, text width=5.2cm] at (3.5, 0) {Vertices in topological order, with the most cities on a route
    from 1 (blue). Each value is one more than the best predecessor; the red route $1 \to 2 \to 4 \to 6$ visits 4 cities.};
\end{tikzpicture}
\caption{Longest path in a DAG by DP in topological order. Parent pointers recover the route.}
\label{fig:dagpath}
\end{figure}

\begin{problem}{Longest Flight Route (CSES)}
Find a route from city 1 to city $n$ that visits as many cities as possible, in a network of one-way flights that contains no
cycles.
\end{problem}

\begin{code}
void solve() {
    int n, m;
    cin >> n >> m;
    vvi adj(n + 1);
    vi indeg(n + 1, 0);
    rep(i, m) {
        int a, b;
        cin >> a >> b;
        adj[a].push_back(b);
        indeg[b]++;
    }
    // Kahn's algorithm gives a topological order of the DAG
    vi order, q;
    rep(v, 1, n + 1) if (!indeg[v]) q.push_back(v);
    while (!q.empty()) {
        int u = q.back();
        q.pop_back();
        order.push_back(u);
        for (int v : adj[u])
            if (--indeg[v] == 0) q.push_back(v);
    }
    // best[v]: most cities on a route 1 -> v (0 = unreachable)
    vi best(n + 1, 0), par(n + 1, 0);
    best[1] = 1;
    for (int u : order) {
        if (!best[u]) continue;
        for (int v : adj[u])
            if (best[u] + 1 > best[v]) best[v] = best[u] + 1, par[v] = u;
    }
    if (!best[n]) {
        print("IMPOSSIBLE");
        return;
    }
    vi route;
    for (int v = n; v; v = par[v]) route.push_back(v);
    reverse(all(route));
    print(sz(route));
    print(route);
}
\end{code}

\begin{summary}[Chapter summary]
\begin{itemize}
  \item All algorithms relax edges; they differ in the order. Parent pointers give the shortest-path tree.
  \item Dijkstra finalises vertices in order of distance and needs non-negative weights; skip outdated queue entries.
  \item Put extra information into the vertex (layered graphs); gather counts in finalisation order.
  \item 0--1 weights: a deque. Negative weights: Bellman--Ford, $n - 1$ rounds, plus one to detect negative cycles.
  \item Longest paths negate weights; they are unbounded only if an improvable cycle is reachable and reaches the target.
  \item Difference constraints are shortest paths from a virtual source; a negative cycle means no solution.
  \item Floyd--Warshall adds allowed intermediate vertices one by one: $\Oh{n^3}$, any weights.
  \item In a DAG, relax in topological order: shortest and longest paths in linear time.
\end{itemize}
\end{summary}

\begin{practice}
\ppgroup{Single source}
\begin{pplist}
  \item Shortest Routes I
  \item Flight Discount
  \item Investigation
  \item Flight Routes
  \item High Score
  \item Cycle Finding
\end{pplist}
\ppgroup{All pairs and DAGs}
\begin{pplist}
  \item Shortest Routes II
  \item Longest Flight Route
  \item Game Routes
\end{pplist}
\end{practice}
